\section{Round-13 Objective}
We pursue \textbf{(C) obstruction/correction}.  The immediately previous round proposed an exact decomposition
\[
A=B\sqcup E,
\]
where
\[
B:=\{n\in\mathbb N:\ n\not\equiv a_i\pmod{n_i}\ \forall i\ge 1\}
\]
is the thresholdless core and \(E\) is the exceptional contribution coming from future congruence classes.  That reduction was promising, but one crucial claim from that round was too strong: it implicitly treated the exceptional set as negligible (indeed, ``automatically logarithmic-density \(0\)'') once the core \(B\) was understood.

This round identifies a fatal obstruction to that strategy.
We prove that the exceptional set can itself have \emph{arbitrarily large logarithmic density}, even in the extreme case where the thresholdless core is empty:
\[
B=\varnothing,
\qquad
\delta_{\log}(A)>1-\varepsilon.
\]
So the previous-round reduction to the thresholdless core alone is false as a route to the full problem.

The most promising direction, given the previous round, is therefore to correct that reduction rigorously and then rule out the whole class of proof strategies that try to solve Problem~25 by proving the thresholdless core is well behaved while declaring the exceptional set negligible.

\section{Round-12 Foundation Used}
We rely on the following specific pieces of the immediately previous round.
\begin{itemize}
\item The definition of the thresholdless core
\[
B:=\{n\in\mathbb N:\ n\not\equiv a_i\pmod{n_i}\ \forall i\ge 1\}.
\]
\item The idea that the difference between \(A\) and \(B\) comes only from the least nonnegative representatives of the residue classes \(a_i\pmod{n_i}\).
\item From Round~11, the corrected summable-reciprocals theorem: if
\[
\sum_{i=1}^{\infty}\frac1{n_i}<\infty,
\]
then \(\delta_{\log}(A)\) exists and equals
\[
\lim_{k\to\infty} d(A^{(k)}),
\qquad
A^{(k)}:=\{n\in\mathbb N:\ n\not\equiv a_i\pmod{n_i}\ \forall 1\le i\le k\}.
\]
\end{itemize}
We do \emph{not} reuse the previous round's stronger characterization of the exceptional set without proof; instead we identify and correct a concrete error in it below.

\section{New Insight / Tool (Round-13)}
The new ingredient has two parts.

\subsection*{(i) Corrected exact decomposition}
If \(r_i\in\{0,1,\dots,n_i-1\}\) denotes the least nonnegative representative of \(a_i\pmod{n_i}\) and
\[
R:=\{r_i:\ i\ge 1\},
\]
then the exact decomposition is
\[
A=B\sqcup (A\cap R).
\]
This is the right formulation.  The stronger-looking previous-round characterization
\[
E=\{r_i:\ r_i\in A^{(i-1)}\}
\]
is false in general.

\subsection*{(ii) Explicit obstruction family}
For any integer \(Q\ge 2\), consider
\[
n_i:=Q^i+i,
\qquad
 a_i\equiv i\pmod{n_i}.
\]
Then every positive integer is the least residue representative of a future congruence class, so the thresholdless core is empty:
\[
B=\varnothing.
\]
Nevertheless, because
\[
\sum_{i\ge 1}\frac1{Q^i+i}<\infty,
\]
Round~11 gives existence of \(\delta_{\log}(A)\), and a simple union bound shows
\[
\delta_{\log}(A)\ge 1-\sum_{i\ge 1}\frac1{Q^i+i}>1-\frac1{Q-1}.
\]
Thus the exceptional set alone can carry logarithmic density arbitrarily close to \(1\).

\section{Attack Plan (Round-13)}
The exact gap left by the previous round is this:
\begin{quote}
Is the thresholdless core \(B\) the only genuinely difficult part, with the exceptional contribution automatically negligible?
\end{quote}
We answer this negatively.

The plan is:
\begin{enumerate}
\item Correct the previous-round exceptional-set formula by proving the exact decomposition
\[
A=B\sqcup (A\cap R).
\]
\item Exhibit a concrete counterexample to the stronger previous-round characterization of \(E\).
\item Produce an explicit infinite family with \(B=\varnothing\).
\item Prove, using the Round~11 summable-reciprocals theorem, that for this family the full set \(A\) (hence the exceptional part) has logarithmic density bounded below by a positive constant, indeed arbitrarily close to \(1\).
\item Conclude that any proof strategy reducing Problem~25 to the thresholdless core alone is fatally obstructed.
\end{enumerate}

\section{Work (Round-13)}
\subsection*{13.1. Corrected exact decomposition}
For each \(i\ge 1\), let \(r_i\in\{0,1,\dots,n_i-1\}\) be the least nonnegative representative of the residue class \(a_i\pmod{n_i}\), and set
\[
R:=\{r_i:\ i\ge 1\}.
\]
Recall
\[
B:=\{n\in\mathbb N:\ n\not\equiv a_i\pmod{n_i}\ \forall i\ge 1\}.
\]

\paragraph{Lemma 13.1 (exact corrected decomposition).}
One has
\begin{equation}
A=B\sqcup (A\cap R).
\tag{13.1}
\end{equation}
In particular,
\begin{equation}
A\setminus B=A\cap R.
\tag{13.2}
\end{equation}

\emph{Proof.}
First, \(B\subseteq A\): if \(n\in B\), then \(n\) avoids every congruence class \(a_i\pmod{n_i}\), hence certainly avoids every \emph{active} congruence class, so \(n\in A\).

Now let \(n\in A\setminus B\).  Since \(n\notin B\), there exists some \(i\) such that
\[
n\equiv a_i\pmod{n_i}.
\]
But \(n\in A\) means that for every active index \(j\) one has \(n\not\equiv a_j\pmod{n_j}\).  Therefore the index \(i\) above cannot be active, so necessarily \(n<n_i\).  Since \(n\) is a positive integer with \(0\le n<n_i\) and is congruent to \(a_i\pmod{n_i}\), it must equal the least residue representative:
\[
n=r_i.
\]
Hence \(n\in A\cap R\).

Conversely, if \(n\in A\cap R\), then \(n=r_i\) for some \(i\), so \(n<n_i\) and
\[
n\equiv a_i\pmod{n_i}.
\]
Thus \(n\notin B\).  Since also \(n\in A\), we have \(n\in A\setminus B\).

So \(A\setminus B=A\cap R\), which proves \eqref{13.2}; together with \(B\subseteq A\), this yields the disjoint union \eqref{13.1}. \hfill \(\square\)

\paragraph{Lemma 13.2 (the previous-round stronger characterization is false).}
The implication
\[
r_i\in A\ \Longrightarrow\ r_i\in A^{(i-1)}
\]
is false in general.

\emph{Proof.}
Take
\[
n_1=5,\ a_1\equiv 1\pmod 5,
\qquad
n_2=7,\ a_2\equiv 1\pmod 7.
\]
Then \(r_2=1\).  Since \(1<n_1\), no constraint is active at \(n=1\), so \(1\in A\).  However,
\[
1\equiv 1\pmod 5,
\]
so \(1\notin A^{(1)}\).  Thus the displayed implication fails. \hfill \(\square\)

\subsection*{13.2. A family with empty thresholdless core but large logarithmic density}
Fix an integer \(Q\ge 3\), and define
\begin{equation}
n_i:=Q^i+i,
\qquad
 a_i:=i \pmod{n_i}
\qquad (i\ge 1).
\tag{13.3}
\end{equation}
Since \(Q^i+i<i+Q^{i+1}=n_{i+1}\), the moduli are strictly increasing.  Also \(0<i<n_i\), so the least nonnegative representative is exactly
\[
r_i=i.
\]
Let \(A\) be the associated set, and let \(B\) be its thresholdless core.

\paragraph{Lemma 13.3 (the thresholdless core is empty).}
For the family \eqref{13.3}, one has
\begin{equation}
B=\varnothing.
\tag{13.4}
\end{equation}
Consequently,
\begin{equation}
A=A\cap R,
\qquad
R=\mathbb N.
\tag{13.5}
\end{equation}

\emph{Proof.}
Let \(n\in\mathbb N\) be arbitrary.  Choosing the index \(i=n\), we have
\[
n<n_n=Q^n+n
\]
and
\[
n\equiv n\equiv a_n\pmod{n_n}.
\]
Therefore \(n\) violates the thresholdless condition defining \(B\).  Since \(n\) was arbitrary, \(B=\varnothing\).

Also \(r_i=i\) for every \(i\), so
\[
R=\{r_i:\ i\ge 1\}=\mathbb N.
\]
Now apply Lemma~13.1:
\[
A=B\sqcup (A\cap R)=\varnothing\sqcup (A\cap\mathbb N)=A.
\]
This proves \eqref{13.5}. \hfill \(\square\)

\paragraph{Lemma 13.4 (uniform positive lower bound for the stage densities).}
Let
\[
A^{(k)}:=\{n\in\mathbb N:\ n\not\equiv a_i\pmod{n_i}\ \forall 1\le i\le k\}
\]
and write
\[
\delta_k:=d(A^{(k)}).
\]
Then for every \(k\ge 1\),
\begin{equation}
\delta_k\ge 1-\sum_{i=1}^{k}\frac1{Q^i+i}
\ge 1-\sum_{i=1}^{\infty}\frac1{Q^i+i}
>1-\frac1{Q-1}.
\tag{13.6}
\end{equation}

\emph{Proof.}
The complement of \(A^{(k)}\) is the union of the \(k\) residue classes
\[
\{n\in\mathbb N:\ n\equiv a_i\pmod{n_i}\},
\qquad 1\le i\le k.
\]
Each such residue class has natural density \(1/n_i=1/(Q^i+i)\).  By the union bound for natural densities,
\[
1-\delta_k=d\!\left(\bigcup_{i=1}^{k}\{n:\ n\equiv a_i\pmod{n_i}\}\right)
\le \sum_{i=1}^{k}\frac1{Q^i+i}.
\]
This proves the first inequality in \eqref{13.6}.  The second is immediate, and the third follows from
\[
\sum_{i=1}^{\infty}\frac1{Q^i+i}<\sum_{i=1}^{\infty}\frac1{Q^i}=\frac1{Q-1}.
\]
\hfill \(\square\)

\paragraph{Theorem 13.5 (the exceptional set alone can have arbitrarily large logarithmic density).}
For the family \eqref{13.3}, the logarithmic density \(\delta_{\log}(A)\) exists and satisfies
\begin{equation}
\delta_{\log}(A)=\lim_{k\to\infty}\delta_k
\ge 1-\sum_{i=1}^{\infty}\frac1{Q^i+i}
>1-\frac1{Q-1}.
\tag{13.7}
\end{equation}
Since \(B=\varnothing\), this means the exceptional set \(A\cap R=A\) has logarithmic density \(>1-1/(Q-1)\).

\emph{Proof.}
Because
\[
\sum_{i=1}^{\infty}\frac1{n_i}=\sum_{i=1}^{\infty}\frac1{Q^i+i}<\infty,
\]
the corrected summable-reciprocals theorem from Round~11 applies.  Therefore \(\delta_{\log}(A)\) exists and equals
\[
\delta:=\lim_{k\to\infty}\delta_k.
\]
Now Lemma~13.4 gives the uniform bound
\[
\delta_k\ge 1-\sum_{i=1}^{\infty}\frac1{Q^i+i}
\qquad (k\ge 1),
\]
so passing to the limit yields
\[
\delta\ge 1-\sum_{i=1}^{\infty}\frac1{Q^i+i}>1-\frac1{Q-1}.
\]
Finally, Lemma~13.3 gives \(B=\varnothing\) and \(A=A\cap R\), so the whole logarithmic density comes from the exceptional part. \hfill \(\square\)

\paragraph{Corollary 13.6 (arbitrarily dense exceptional-only examples).}
For every \(\varepsilon>0\), there exists an instance of Problem~25 for which
\begin{equation}
B=\varnothing,
\qquad
\delta_{\log}(A)>1-\varepsilon.
\tag{13.8}
\end{equation}
Indeed, it suffices to choose any integer \(Q>1+1/\varepsilon\) and use the family \eqref{13.3}.

\emph{Proof.}
Choose \(Q\) so large that \(1/(Q-1)<\varepsilon\).  Then Theorem~13.5 gives
\[
\delta_{\log}(A)>1-\frac1{Q-1}>1-\varepsilon.
\]
The identity \(B=\varnothing\) comes from Lemma~13.3. \hfill \(\square\)

\paragraph{Minimal corrected statement.}
The strongest exact reduction justified by the previous round is
\[
A=B\sqcup (A\cap R).
\]
However, there is \emph{no} general negligibility statement for \(A\cap R\): by Corollary~13.6, this exceptional part can already have logarithmic density arbitrarily close to \(1\) while \(B=\varnothing\).

So any successful resolution of Problem~25 must control the exceptional part as a genuine second component; controlling the thresholdless core alone is insufficient.

\section{Adversarial Verification}
We try to break the new work.
\begin{itemize}
\item \textbf{Correctness of the decomposition.}  The proof of Lemma~13.1 uses only the definition of \(A\) and the fact that if \(n<n_i\) and \(n\equiv a_i\pmod{n_i}\), then \(n\) must equal the least nonnegative representative \(r_i\).  No hidden assumptions are used.

\item \textbf{The previous-round stronger formula really fails.}  Lemma~13.2 gives a concrete two-modulus counterexample, so we are not merely expressing doubt; we identify and prove an actual error.

\item \textbf{Why \(B=\varnothing\) in the obstruction family.}  For each \(n\), the single index \(i=n\) already excludes \(n\) from the thresholdless core because \(n\equiv a_n\pmod{n_n}\).  This is exact and requires no asymptotic argument.

\item \textbf{Why \(R=\mathbb N\).}  In the family \eqref{13.3}, the least residue representative is \(r_i=i\), so every positive integer appears as one of the representatives.  Hence \(A=A\cap R\) once \(B=\varnothing\).

\item \textbf{Use of Round~11 is legitimate.}  The only imported theorem is the corrected summable-reciprocals result, and its hypothesis
\(
\sum_i 1/n_i<\infty
\)
holds because \(n_i=Q^i+i\).

\item \textbf{Lower bound on stage densities.}  Lemma~13.4 uses only the trivial union bound on the complement of \(A^{(k)}\).  It does not rely on independence or coprimality.

\item \textbf{What this does \emph{not} prove.}  We do \emph{not} claim that the exceptional set can realize every possible logarithmic density or every oscillatory behaviour.  What is proved is already enough to kill the previous reduction: the exceptional part can have logarithmic density arbitrarily close to \(1\) even when the core is empty.
\end{itemize}

\section{Final}
\textbf{UNRESOLVED (BUT STRICTLY ADVANCED).}

This round proves a fatal obstruction to the immediately previous reduction.

\begin{enumerate}
\item The exact decomposition is
\[
A=B\sqcup (A\cap R),
\]
with \(R\) the set of least residue representatives of the classes \(a_i\pmod{n_i}\).  The stronger characterization used in the previous round is false.

\item More importantly, the exceptional part is \emph{not} negligible in any general sense.  For every \(\varepsilon>0\) there exists an explicit family with
\[
B=\varnothing,
\qquad
\delta_{\log}(A)>1-\varepsilon.
\]
Equivalently, the exceptional part alone can carry logarithmic density arbitrarily close to \(1\).
\end{enumerate}

So the thresholdless-core approach is not a valid reduction of Problem~25.  The minimal corrected statement is that one must understand both pieces in the exact decomposition
\[
A=B\sqcup (A\cap R),
\]
and the second piece can already be essentially as large as possible.

\section{Completion Estimate}
\textbf{COMPLETION: 98\%}.

\section{References}
\begin{itemize}
\item T. F. Bloom, \emph{Erd\H{o}s Problem \#25}, official problem page, \texttt{https://www.erdosproblems.com/25}, accessed 2026-03-09.
\item No external mathematical references were used in the proofs of this round.
\end{itemize}
